\documentclass[11pt]{article}
\usepackage[margin=1in,top=.9in,bottom=.9in]{geometry}
\usepackage{amsmath,amssymb,amsthm,hyperref}
\hypersetup{colorlinks=true,linkcolor=blue,urlcolor=blue,citecolor=blue}
\setlength{\emergencystretch}{2em}
\newtheorem{theorem}{Theorem}
\newtheorem{lemma}[theorem]{Lemma}
\title{Four interior roots in Yuster's entropy polynomial}
\author{Research proof and formalization packet}
\date{October 26, 2026}
\begin{document}
\maketitle
\begin{abstract}
We disprove the polynomial-root conjecture in Yuster's \emph{Almost
$k$-union closed set systems}: at $k=48$, its canonically parameterized
polynomial has at least four distinct roots in $(0,1)$. Five exact
rational sign checks, together with rational bounds on the canonical
algebraic parameter, provide a short certificate. Lean verifies the
literal source coefficients, the parameter, every sign, the four real
roots, and the resulting root count with multiplicity. No claim is made
about the smallest counterexample or the exact total number of roots.
The underlying entropy inequality is a different statement and has
already been proved in subsequent work.
\end{abstract}

\section{The exact source polynomial}
Fix an integer $k\geq2$. Following Yuster \cite[Section~3]{Yuster}, let
$\psi_k$ be the unique solution in $[0,1]$ of
\[
 (1-\psi_k)^k=\psi_k,
 \qquad \phi_k=1-\psi_k,\qquad \alpha_k=\phi_k^{k-1}.
\]
The coefficients in Lemma~3.5 of that source are determined by
\begin{align*}
 C(k,0,0)&=\frac1{(k-1)!},\\
 C(k,0,j)&=0\quad(j>0),\qquad C(k,t,0)=0\quad(t>0),\\
 C(k,t,j)&=(kj-t+1)C(k,t-1,j)+kC(k,t-1,j-1)
       \quad(t,j>0).
\end{align*}
For concise display, set $G_k(x)=\sum_{i=0}^{k-1}x^i$. Corollary~3.7
defines $p_k(x)=\alpha_k\rho_k(x)-\sigma_k(x)$, where
\begin{align}
 \rho_k(x)&=\sum_{j=0}^{k-1}(-1)^j j! C(k,k+1,j+2)\notag\\[-2mm]
 &\hspace{12mm}\cdot\left[(x^k-1)^{k-j-1}x^{kj+k}-(x^k-1)^k\right],\label{eq:rho}\\
 \sigma_k(x)&=\sum_{j=0}^{k-1}(-1)^j\binom{k+1}{j+2}\notag\\[-2mm]
 &\hspace{12mm}\cdot\left[x^{j+1}(x-1)^{k-j-1}G_k(x)^k-(x^k-1)^k\right].\label{eq:sigma}
\end{align}
Conjecture~3.9 asserts that $p_k$ has at most two real roots in $(0,1)$,
counted with multiplicity, for every $k\geq2$. This is the assertion
recorded as MathDB~357503.

\begin{theorem}\label{thm:counterexample}
The polynomial $p_{48}$ defined above has at least four distinct roots
in $(0,1)$. In particular its number of interior roots, counted with
multiplicity, is at least four, and Yuster's Conjecture~3.9 is false.
\end{theorem}

\section{A rational sign certificate}
Let
\[
 a=\frac{61}{1000},\qquad b=\frac{613}{10000}.
\]
First, the source parameter satisfies $a<\alpha_{48}<b$. Indeed
$\phi_k+\phi_k^k=1$ and $\phi_k>0$, so
$\phi_k(1+\alpha_k)=1$ and hence
\begin{equation}\label{eq:alpha}
 \alpha_k(1+\alpha_k)^{k-1}=1.
\end{equation}
Exact rational arithmetic gives
\[
 a(1+a)^{47}<1<b(1+b)^{47}.
\]
The function $u\mapsto u(1+u)^{47}$ is strictly increasing on
$[0,\infty)$, proving the asserted bounds from \eqref{eq:alpha}.
Conversely the certificate constructs the unique source $\psi_{48}$
in $[0,1]$; it does not assume an arbitrary coefficient in this bracket.

The following table records exact rational signs, using the literal
finite sums \eqref{eq:rho}--\eqref{eq:sigma}. A minus or plus denotes a
strict inequality, not a floating-point approximation.
\begin{center}
\begin{tabular}{c|c|c|c}
 $x$ & $\rho_{48}(x)$ & endpoint expression & its sign\\\hline
 $1/100$ & $<0$ & $a\rho_{48}(x)-\sigma_{48}(x)$ & $-$\\
 $1/10$  & $<0$ & $b\rho_{48}(x)-\sigma_{48}(x)$ & $+$\\
 $1/2$   & $<0$ & $a\rho_{48}(x)-\sigma_{48}(x)$ & $-$\\
 $3/5$   & $<0$ & $b\rho_{48}(x)-\sigma_{48}(x)$ & $+$\\
 $4/5$   & $<0$ & $a\rho_{48}(x)-\sigma_{48}(x)$ & $-$
\end{tabular}
\end{center}
For reproducibility, one can calculate all entries using integer rows
$B_{t,j}=47!C(48,t,j)$, starting from $B_{0,0}=1$ and applying
\[
 B_{t+1,j}=(48j-t)B_{t,j}+48B_{t,j-1}.
\]
Undefined row entries are zero. Divide by $47!$ when substituting in
\eqref{eq:rho}; all ensuing evaluations are rational. The packet
contains the full rational values, and Lean independently checks the
inequalities by kernel reduction of exact rational arithmetic.

Since $\rho_{48}(x)<0$ at every listed point, increasing the parameter
decreases $\alpha\rho_{48}(x)-\sigma_{48}(x)$. Thus the bounds
$a<\alpha_{48}<b$ transfer the five signs in the table to $p_{48}$:
\[
 p_{48}(1/100)<0, \ p_{48}(1/10)>0, \ p_{48}(1/2)<0, \
 p_{48}(3/5)>0, \ p_{48}(4/5)<0.
\]
Polynomial continuity and the intermediate value theorem give a root
in each of the four open intervals
\[
 (1/100,1/10),\quad(1/10,1/2),\quad(1/2,3/5),\quad(3/5,4/5).
\]
The intervals are pairwise disjoint and contained in $(0,1)$, so these
four roots are distinct. The polynomial is nonzero by its first strict
sign. Each of the four roots therefore occurs with positive
multiplicity in its root multiset. This proves
Theorem~\ref{thm:counterexample}.

\section{Formal coverage and reproducibility}
The seven modules \texttt{YusterCoefficients},
\texttt{YusterPolynomial}, \texttt{YusterArithmetic},
\texttt{YusterAlpha}, \texttt{YusterSigns}, \texttt{FourRoots}, and
\texttt{YusterRootCounterexample} contain 53 local declarations.
The source recurrence, including its factorial normalization, is
proved for the implemented exact coefficients. The real polynomial is
literally \eqref{eq:rho}--\eqref{eq:sigma}; no auxiliary proxy
polynomial or unproved derivative identity is assumed. The source
parameter is constructed and shown unique in its specified interval.
Rational comparisons are checked with ordinary kernel evaluation
(\texttt{decide +kernel}), not a native computation axiom. Real sign
transfer, intermediate roots, distinctness, and the multiset count are
also certified.

In namespace \texttt{YusterCounterexample}, the main theorem
\path{source_rootcount_ge_four}
counts the roots of the actual polynomial with its canonical
parameter. The further endpoint
\path{source_conjecture_357503_false}
negates the source's universal statement with its defining equation
for $\psi_k$. No sign, canonical-parameter, root-existence, or
source-identity assumption is left to the user.

A fresh isolated replay compiled all declarations with warnings as
errors under Lean~4.34.1 and the pinned mathlib dependencies. The
compiler and mathlib commits, respectively, are
\begin{center}\small\ttfamily
5045d0056413266e57c625dcd7c365b10e377c52\\
d13f23b723b8a846827a245b89c10fc7d3f11612.
\end{center}
Axiom inspection found only \texttt{propext},
\texttt{Classical.choice}, and \texttt{Quot.sound}. The packet includes
the replay script, exact dependency pins, declaration types and axiom
inventory, and independent review. No proof holes or added axioms are
used.

\section{Relation to subsequent work and novelty limits}
This counterexample concerns the proposed polynomial root bound.
It does not refute an entropy inequality or a union-closed-set theorem.
Ho \cite{Ho} proves the generalized entropy inequality by another
route; that result is credited and is not claimed here. Wakhare
\cite[Conjecture~2]{Wakhare} proposes a closely related polynomial
root bound. The packet claims one outcome, and does not present a
separate Lean refutation of Wakhare's formulation.

The source review on October~26, 2026 inspected Yuster's complete
primary manuscript, including its actual coefficient formulas and
Conjecture~3.9. Only revision v1 was listed. MathDB's references field
contains an inconsistent ``solved'' description quoting an already
resolved $k=2$ entropy constant from the source introduction; it is not
a resolution of the root conjecture on page~9. No posted solution of
the root assertion was present in the reviewed record.

Targeted primary-source searches and independent review located no
earlier root-count counterexample. The complete later
Yuster--Yashfe manuscript was not located, and the published Wakhare
body was not separately compared with its latest arXiv v2. These are
explicit limits on the prior-art assessment. The result is offered
with a documented, bounded novelty check, not an exhaustive worldwide
priority claim. Neither the exact total number of roots nor the
minimality of $k=48$ is asserted.

\begin{thebibliography}{9}
\bibitem{Yuster}
R.~Yuster, \emph{Almost $k$-union closed set systems},
\href{https://arxiv.org/abs/2302.12276v1}{arXiv:2302.12276v1}
(2023), Lemma~3.5, Corollary~3.7, Conjecture~3.9.
\bibitem{Wakhare}
T.~Wakhare, \emph{Iterated Entropy Derivatives and Binary Entropy
Inequalities}, Journal of Approximation Theory \textbf{307} (2025),
106143; primary manuscript
\href{https://arxiv.org/abs/2312.14743v2}{arXiv:2312.14743v2},
Conjecture~2.
\bibitem{Ho}
B.S.~Ho, \emph{A generalization of Boppana's entropy inequality},
\href{https://arxiv.org/abs/2601.19327v1}{arXiv:2601.19327v1}
(2026).
\end{thebibliography}
\end{document}
